\iffalse Even though the world abounds with many rich and diverse
dialects, each possessing distinct features and characteristics, many
of these dialects still lack the resources and support required for
their speakers to access contemporary technologies in their language
\cite{joshi-etal-2020-state}.  Therefore, efforts to create annotated
corpora, develop language models, or establish dictionaries and
grammars for low-resource dialects are critical for preserving and
promoting these dynamic languages, which capture the community's
unique cultures, histories, and experiences. For example, the
Masakhane community aims to strengthen NLP research for African
languages with important initiatives such as MasakhaNER
\cite{adelani-etal-2021-masakhaner}. By working on resources for
low-resource dialects, we contribute to their preservation and
promotion.  
\else 

Despite the abundance of rich and diverse dialects
worldwide, each possessing distinctive features and characteristics,
many of these dialects still lack the necessary resources and support
to enable their speakers to access modern technologies in their own
language \cite{joshi-etal-2020-state}. Therefore, it is imperative
to undertake endeavors aimed at creating annotated corpora, developing
language models, and establishing dictionaries and grammars for
low-resource dialects. These efforts are crucial for the preservation
and advancement of these dynamic languages, which encapsulate unique
cultures, histories, and experiences within their respective
communities. 

One notable example of such an effort is the Masakhane community, which is
dedicated to enhancing natural language processing (NLP) research for
African languages through significant initiatives such as MasakhaNER
\cite{adelani-etal-2021-masakhaner}. Similar efforts are ongoing for
Indonesian languages \cite{cahyawijaya2022nusacrowd}.




\iffalse 

Another long standing, orthogonal in a way, established initiative
that initially aimed at offering a common set of syntactic guidelines
for a small set of languages, the Universal Dependencies project
\cite{nivre-etal-2020-universal}, turned out to be the recipient of
many low-resources languages treebank initiatives that followed and
often extended these initial guidelines to cope with their respective
idiosyncrasies.

\else 
In addition, a long-standing and somewhat unrelated initiative
known as the Universal Dependencies project
\cite{nivre-etal-2020-universal} originally aimed to provide a
standardized set of syntactic guidelines for a limited number of
languages turned out to become the recipient of numerous
treebank initiatives for low-resource languages. These initiatives not
only adopted the initial guidelines but also expanded upon them to
accommodate the unique idiosyncrasies of each language.  \fi

In this work, we aim to enhance a pre-existing multi-view
treebank  devoted to a very low-resource language, namely the North-African Arabic dialect written in
Latin script, collected from Algerian sources and denoted as the
Narabizi treebank, the first
 available for this dialect, where Arabizi refers to both the practice of
 writing Arabic using the Latin alphabet and \textit{N} for the North
 African dialect  \cite{seddah-etal-2020-building}. 
Made of noisy user-generated content that exhibits a high level of
language variability, its annotations faced many challenges as
described by the authors and contained remaining errors \cite{touileb-barnes-2021-interplay}. 

Our work builds on previous efforts to annotate and standardize
treebank annotations for low-resource languages to enhance the quality
and consistency of linguistic resources
\cite{schluter2007preparing,sade-etal-2018-hebrew,turk-etal-2019-turkish,zariquiey-etal-2022-building}.

 Following previous research, we consider the impact of refining
 annotation schemes on downstream tasks.
 \citet{mille-etal-2012-granularity} examine how much a treebank's
 performance relies on its annotation scheme and whether employing a
 more linguistically rich scheme would decrease performance. Their
 findings indicate that using a fine-grained annotation for training a
 parser does not necessarily improve performance when parsing with a
 coarse-grained tagset. This observation is relevant to our study as
 we expect refining the treebank could enhance the parsing
 performance even though the inherent variability of this language,
 which, tied to its  small size treebank, could bring a negative impact on such
 enhancements.  

 On the other hand, the experiments conducted by
 \citet{schluter2007preparing} demonstrate that using a cleaner and
 more coherent treebank yields superior results compared to a treebank
 with a training set five times larger. This observation highlights
 the significance of high-quality dataset annotations, particularly
 for smaller datasets. This understanding primarily drives the goal of
 improving the NArabizi treebank's annotations.

 In this context, we propose a heavily revised version of NArabizi
 treebank \cite{seddah-etal-2020-building} that includes two novel
 annotation layers for Named Entity Recognition (NER) and offensive
 language detection. One of the goals of this work is also to study the
 impact of non-gold tokenization on NER, a scenario almost never investigated
by the community \cite{bareket-tsarfaty-2021-neural}.
Our primary contributions are as follows:

\begin{itemize}
    \item Using error mining tools, we release a new corrected version
      of the treebank, which leads to improved downstream task performance.
    \item We show that corrections made to a small size treebank of a
      highly variable language favorably impacts the performance of NLP models trained on it.
    \item We augment the treebank by adding NER annotations and offensive language detection, expanding its applicability in various NLP tasks.
    \item We homogenize tokenization across the dataset, analyze the
      impact of proper tokenization on UD tasks and NER and conduct a
      realistic evaluation on predicted tokenization, including 
      NER evaluation.
\end{itemize}

The enhanced version of the Narabizi Treebank is freely available.\footnote{\url{https://gitlab.inria.fr/ariabi/release-narabizi-treebank}}